%
% v2-architecture.tex
%
% Lux v2 Architecture — bones and interface sketches.
% Written at hypothesis stage; details will evolve.
%
% Compile: pdflatex v2-architecture.tex
%

\documentclass[a4paper,10pt]{article}
\usepackage[margin=1in]{geometry}
\usepackage{booktabs}
\usepackage{listings}
\usepackage{xcolor}
\usepackage{enumitem}
\usepackage{longtable}
\usepackage{tabularx}
\usepackage[colorlinks=true,linkcolor=blue,citecolor=blue,urlcolor=blue]{hyperref}

\lstset{
  basicstyle=\ttfamily\small,
  keywordstyle=\color{blue!80!black},
  commentstyle=\color{green!50!black},
  stringstyle=\color{red!60!black},
  breaklines=true,
  frame=single,
  framerule=0.3pt,
  rulecolor=\color{gray!50},
  backgroundcolor=\color{gray!5},
  columns=fullflexible,
}

\newcommand{\api}[1]{\texttt{#1}}
\newcommand{\field}[1]{\texttt{#1}}
\newcommand{\ext}[1]{\texttt{#1}}

\begin{document}

\title{Lux v2: Architecture Sketch}
\author{Agent IDE for Python-Native Visual Output}
\date{April 2026 --- Hypothesis Stage}
\hypersetup{
  pdftitle={Lux v2: Architecture Sketch},
  pdfauthor={Punt Labs},
  pdfsubject={Architecture},
}
\maketitle

\tableofcontents
\newpage

% ===================================================================
\section{Overview}
% ===================================================================

Lux~v2 is a Python-native display compositor for AI agents.  It
provides a persistent local service that agents connect to, renders
their content, and grows new capabilities on demand through
LLM-authored extensions.

\subsection{Design Principles}

\begin{enumerate}[label=\textbf{P\arabic*}.]
  \item \textbf{Everything is an extension.}  The core provides
    transport, dispatch, registry, frame management, render loop, and
    safety wrapper.  Every user-visible element kind, applet, theme,
    and service is an extension.
  \item \textbf{Closed core, open extension surface.}  The core is
    small and stable.  Extensions are the growth vector.  Open-Closed
    Principle applied at the system level.
  \item \textbf{Python is the substrate.}  Not web.  Extensions use
    Python's package ecosystem directly: matplotlib, plotly, PIL,
    pyte, Qt, tkinter.  No transpilation to JavaScript.
  \item \textbf{Ollama-shaped service.}  The primary interface is HTTP
    endpoints on localhost.  Any agent, any language, any machine
    (with TCP transport) can drive the display.  MCP is an adapter
    layer.
  \item \textbf{Self-extending by default.}  When an agent needs a
    visualization that does not exist, it can author one.  The
    extension persists for the user and optionally for the community.
  \item \textbf{Bidirectional.}  The display is not passive.  Agents
    write to Lux; Lux writes back (events, menu clicks, interaction
    messages, screenshots).  The display can drive agent behavior.
\end{enumerate}

\subsection{Relationship to v1}

Lux~v1 is a fixed 27-element vocabulary with a hardcoded ImGui
dispatcher.  v2 preserves the wire protocol shape, the frame system,
the ImGui render loop, and the MCP adapter --- but refactors element
dispatch from hardcoded to registry-based and adds the extension
system, service API, and in-display agent.

% ===================================================================
\section{System Architecture}
% ===================================================================

\subsection{Component Map}

\begin{center}
\small
\begin{tabularx}{\linewidth}{@{}l>{\raggedright\arraybackslash}X@{}}
\toprule
\textbf{Component} & \textbf{Responsibility} \\
\midrule
\textbf{Core} & \\
\quad Transport     & Unix socket (local) + TCP/TLS (remote).
  Length-prefixed JSON framing (inherited from v1). \\
\quad Protocol      & Message types, element schemas, serialization.
  Extended with extension-declared element kinds. \\
\quad Registry      & Runtime dispatch table: element kind $\to$
  renderer.  Populated by extension loader at startup,
  extensible at runtime. \\
\quad Loader        & Scans extension directories, validates manifests,
  imports renderers, registers with the registry.
  Crash-isolates each extension. \\
\quad Frame Manager & Frame lifecycle, shared ownership, orphan model,
  dock bar, focus coordination.  Inherited from v1. \\
\quad Render Loop   & ImGui event loop via \texttt{immapp.run()}.
  Polls sockets, dispatches to registry, flushes events.
  Single-threaded. \\
\quad Safety Wrapper & Exception isolation per extension render call.
  Crash $\to$ error placeholder.  Consent dialog on install. \\
\midrule
\textbf{Service Layer} & \\
\quad HTTP Server   & Localhost HTTP endpoints (ollama-shaped).
  Runs in a thread; posts to main-thread queue. \\
\quad /help         & Auto-documenting endpoint.  Returns installed
  element kinds, services, conventions, permission mode. \\
\quad MCP Adapter   & Thin translation: MCP tool calls $\to$ service
  API calls.  Preserves Claude Code integration. \\
\midrule
\textbf{Extensions} & \\
\quad Starter Pack  & Bundled extensions: text, button, table, plot,
  markdown, image, draw, etc.  Ship with install. \\
\quad User Local    & \texttt{\char`\~/.lux/extensions/}.  Persists
  across sessions. \\
\quad Project Local & \texttt{.lux/extensions/}.  Scoped to repo. \\
\quad Community     & Registry of published extensions.  Installed
  via \texttt{lux ext install}. \\
\midrule
\textbf{Agent SDK} (optional) & \\
\quad LuxExchange   & Tool-use loop (Anthropic Python SDK).  Drives
  the in-display agent. \\
\quad Commands      & Leaf + composite tool commands operating on the
  display.  Transactional composites with rollback. \\
\quad Permission    & Five-mode policy: Plan, Default, Accept-edits,
  Bypass, DontAsk. \\
\bottomrule
\end{tabularx}
\normalsize
\end{center}

\subsection{Data Flow}

\begin{verbatim}
External Agent (Claude Code, CLI, remote)
  |
  |  HTTP or MCP
  v
Service Layer (thread)
  |
  |  Queue (thread-safe)
  v
Main Thread: Render Loop
  |
  +-- Registry lookup: element kind -> extension renderer
  |
  +-- Extension renderer: imgui | texture | window
  |
  +-- Frame Manager: position, focus, dock
  |
  +-- Event flush: InteractionMessages -> all clients
\end{verbatim}

% ===================================================================
\section{Extension System}
% ===================================================================

\subsection{Extension Manifest}

Every extension declares itself in a manifest file
(\texttt{extension.py} or \texttt{extension.toml}) at the extension
root:

\begin{lstlisting}[language=Python]
# ~/.lux/extensions/sankey/extension.py

LUX_EXTENSION = {
    "name": "sankey",
    "version": "1.0.0",
    "author": "maya-torres",
    "description": "Sankey flow diagram",
    "element_kind": "sankey",
    "render_mode": "texture",     # imgui | texture | window
    "schema": {
        "nodes": {"type": "list", "required": True},
        "edges": {"type": "list", "required": True},
        "title": {"type": "str",  "required": False},
    },
    "dependencies": ["plotly>=5.0"],
    "lux_version": ">=2.0",
}

def render(element, ctx):
    """Render a sankey diagram.

    Args:
        element: Validated element dict (matches schema).
        ctx: RenderContext with mode-specific capabilities.
    """
    import plotly.graph_objects as go
    fig = go.Figure(go.Sankey(
        node=dict(label=[n["label"] for n in element["nodes"]]),
        link=dict(
            source=[e["source"] for e in element["edges"]],
            target=[e["target"] for e in element["edges"]],
            value=[e["value"] for e in element["edges"]],
        ),
    ))
    fig.update_layout(title=element.get("title", ""))
    ctx.render_figure(fig)  # texture mode: fig -> PNG -> texture
\end{lstlisting}

\subsection{Three Rendering Modes}

Each extension declares one of three rendering modes.  The mode
determines the \texttt{RenderContext} capabilities:

\begin{center}
\begin{tabularx}{\linewidth}{@{}lp{3cm}>{\raggedright\arraybackslash}X@{}}
\toprule
\textbf{Mode} & \textbf{ctx Provides} & \textbf{When to Use} \\
\midrule
\ext{imgui} &
  \texttt{imgui} module,
  \texttt{draw\_list},
  \texttt{widget\_state},
  \texttt{send\_event()} &
  Native ImGui widgets, custom drawing, interactive controls.
  Preferred mode --- single-window UX, full interactivity. \\
\ext{texture} &
  \texttt{render\_image(pil\_img)},
  \texttt{render\_figure(mpl\_fig)},
  \texttt{render\_bytes(png)} &
  Libraries that render to bitmaps: matplotlib, PIL, plotly
  (static export), any numpy$\to$pixels pipeline.  Displayed
  as ImGui image in the Lux window.  Limited interactivity
  (click forwarding, no native zoom/pan). \\
\ext{window} &
  \texttt{create\_window(title)},
  \texttt{on\_close(callback)},
  \texttt{send\_event()} &
  Libraries that require their own OS window: Qt, tkinter,
  Jupyter kernels, Pharo (via Postern).  Lux tracks lifecycle
  and coordinates focus but does not composite pixels.
  Escape hatch for anything ImGui cannot host. \\
\bottomrule
\end{tabularx}
\end{center}

\subsection{Extension Lifecycle}

\begin{enumerate}
  \item \textbf{Discovery.}  Loader scans directories in order:
    bundled starter pack, project-local
    (\texttt{.lux/extensions/}), user-local
    (\texttt{\char`\~/.lux/extensions/}).
  \item \textbf{Validation.}  Manifest parsed.  Schema checked.
    Dependencies verified via \texttt{importlib.metadata}.
    Incompatible \field{lux\_version} $\to$ skipped with warning.
  \item \textbf{Registration.}  Element kind added to registry.
    Renderer function stored in dispatch table.
  \item \textbf{Rendering.}  When a scene contains an element with
    \field{kind} matching this extension, the registry dispatches
    to the extension's \texttt{render()} function with a
    mode-appropriate \texttt{RenderContext}.
  \item \textbf{Crash isolation.}  Every \texttt{render()} call is
    wrapped in \texttt{try/except}.  Exception $\to$ error
    placeholder rendered in place; core keeps running.  Repeated
    crashes $\to$ extension disabled for the session.
  \item \textbf{Hot install.}  \api{lux ext install} writes to
    user-local directory and signals the display process to re-scan.
    New extension available next frame.  No restart required.
\end{enumerate}

\subsection{Extension Registry API}

The registry is a Python dict at runtime:

\begin{lstlisting}[language=Python]
_registry: dict[str, ExtensionEntry] = {}

@dataclass
class ExtensionEntry:
    name: str
    kind: str                     # element kind
    render_mode: str              # "imgui" | "texture" | "window"
    schema: dict                  # field validation
    renderer: Callable            # render(element, ctx)
    metadata: dict                # author, version, description
    crash_count: int = 0          # for auto-disable
\end{lstlisting}

% ===================================================================
\section{Service API}
% ===================================================================

The primary interface to Lux~v2 is HTTP on localhost (default port
8423).  The shape follows Ollama: simple REST endpoints, JSON
payloads, streaming where appropriate.

\subsection{Endpoints}

\begin{center}
\begin{tabularx}{\linewidth}{@{}llX@{}}
\toprule
\textbf{Method} & \textbf{Path} & \textbf{Purpose} \\
\midrule
POST & \api{/api/show}      & Render a scene (elements, frame
  targeting). \\
POST & \api{/api/update}    & Patch elements by ID. \\
GET  & \api{/api/events}    & Stream interaction events (SSE). \\
GET  & \api{/api/frames}    & Introspect current frame state. \\
POST & \api{/api/clear}     & Clear all scenes and frames. \\
GET  & \api{/api/ping}      & Health check with round-trip time. \\
\midrule
GET  & \api{/api/extensions}          & List installed extensions. \\
POST & \api{/api/extensions/install}  & Install from path or URL. \\
POST & \api{/api/extensions/author}   & LLM-authored extension
  (requires in-display agent). \\
DELETE & \api{/api/extensions/\{name\}} & Remove an extension. \\
\midrule
GET  & \api{/api/capabilities} & Installed kinds, services,
  render modes. \\
GET  & \api{/help}           & Auto-documenting: conventions,
  schemas, permission mode.  Same contract as
  Postern's \api{/help}. \\
\midrule
POST & \api{/api/screenshot} & Capture display as PNG.  Returns
  base64 or writes to path. \\
POST & \api{/api/agent/run}  & Submit a prompt to the in-display
  agent (requires Agent SDK). \\
\bottomrule
\end{tabularx}
\end{center}

\subsection{The \api{/help} Contract}

\api{/help} returns runtime truth from the live registry, not
static documentation.  An agent that has never seen this Lux
instance reads \api{/help} first and learns:

\begin{itemize}[nosep]
  \item What element kinds are installed right now
  \item What schemas each kind expects
  \item What rendering modes are active
  \item What permission mode is configured
  \item What extensions are available but not loaded (disabled,
    missing deps)
  \item Conventions and safety rules
\end{itemize}

This is the same pattern as Postern's \api{/help} endpoint:
the live system self-documents.

\subsection{MCP Adapter}

The MCP server (\texttt{lux serve}) becomes a thin translation
layer:

\begin{center}
\begin{tabular}{@{}ll@{}}
\toprule
\textbf{MCP Tool} & \textbf{Service Call} \\
\midrule
\texttt{show()}          & \api{POST /api/show} \\
\texttt{update()}        & \api{POST /api/update} \\
\texttt{recv()}          & \api{GET /api/events} (single event) \\
\texttt{clear()}         & \api{POST /api/clear} \\
\texttt{ping()}          & \api{GET /api/ping} \\
\texttt{display\_mode()} & \api{GET/POST /api/config/display} \\
\texttt{show\_table()}   & \api{POST /api/show} (convenience) \\
\texttt{show\_dashboard()} & \api{POST /api/show} (convenience) \\
\texttt{show\_diagram()} & \api{POST /api/show} (convenience) \\
\bottomrule
\end{tabular}
\end{center}

Existing Claude Code integration is preserved.  Agents that want
the richer surface (extension management, streaming events,
agent interaction) use the service API directly.

% ===================================================================
\section{In-Display Agent SDK}
% ===================================================================

Optional component.  Same architecture as the Claude Agent SDK for
Smalltalk: a tool-use loop where tools operate on the live display
process.

\subsection{LuxExchange}

\begin{lstlisting}[language=Python]
exchange = LuxExchange(
    client=anthropic.Client(),
    commands=[
        ShowSceneCommand(),
        InstallExtensionCommand(),
        AuthorExtensionCommand(),
        IntrospectFramesCommand(),
        ScreenshotCommand(),
        RunPythonCommand(),       # consent-gated
    ],
    permission_policy=DefaultPolicy(),
    max_iterations=20,
)
exchange.run("Render this data as a sankey diagram")
\end{lstlisting}

\subsection{Tool Commands}

Following the Gang-of-Four Composite pattern from the Smalltalk SDK:

\begin{center}
\begin{tabularx}{\linewidth}{@{}lll>{\raggedright\arraybackslash}X@{}}
\toprule
\textbf{Command} & \textbf{Type} & \textbf{Mode} & \textbf{Purpose} \\
\midrule
ShowScene        & Leaf & Safe   & Render a scene. \\
UpdateScene      & Leaf & Safe   & Patch elements. \\
ClearScene       & Leaf & Safe   & Clear frames. \\
IntrospectFrames & Leaf & Safe   & Dump current state. \\
ListExtensions   & Leaf & Safe   & Registry query. \\
ListCapabilities & Leaf & Safe   & What kinds are available. \\
Screenshot       & Leaf & Safe   & Capture display as image. \\
InstallExtension & Leaf & Gated  & Add extension (consent). \\
RemoveExtension  & Leaf & Gated  & Remove extension. \\
AuthorExtension  & Leaf & Gated  & Generate via LLM. \\
RunPython        & Leaf & Gated  & Exec in display namespace. \\
\midrule
AuthorAndInstall & Composite & Gated & Author + install + verify.
  Rollback on failure. \\
RenderWithExt    & Composite & Gated & Install if needed + render.
  Rollback on render failure. \\
\bottomrule
\end{tabularx}
\end{center}

\subsection{Permission Policy}

Five modes, same as the Smalltalk SDK:

\begin{center}
\begin{tabular}{@{}lp{8cm}@{}}
\toprule
\textbf{Mode} & \textbf{Behavior} \\
\midrule
Plan        & Show plan, do not execute.  For review. \\
Default     & Ask per gated call.  Consent dialog. \\
Accept-edits & Auto-approve install/author.  Ask for
  RunPython. \\
Bypass      & Approve everything.  For trusted pipelines. \\
DontAsk     & Deny all gated calls.  Safe mode. \\
\bottomrule
\end{tabular}
\end{center}

\subsection{Screenshot-Driven Iteration}

The in-display agent can take a screenshot of its own output and
feed it back to the LLM for visual QA:

\begin{enumerate}[nosep]
  \item Agent authors an extension and renders content.
  \item Agent calls \texttt{ScreenshotCommand}.
  \item Screenshot sent to Claude as an image content block.
  \item Claude evaluates: ``the columns are misaligned.''
  \item Agent modifies the extension and re-renders.
  \item Repeat until visually correct.
\end{enumerate}

This technique was proven at scale during the Postern/Pharo work,
where agents iteratively refined Morphic UIs using screenshots
of the running image.

% ===================================================================
\section{Community Registry}
% ===================================================================

\subsection{Extension Distribution}

Extensions are distributed as directories containing:

\begin{itemize}[nosep]
  \item \texttt{extension.py} or \texttt{extension.toml} --- manifest
  \item \texttt{renderer.py} --- render function (or inline in
    manifest)
  \item \texttt{requirements.txt} --- Python dependencies
  \item \texttt{tests/} --- test suite
  \item \texttt{README.md} --- usage documentation
\end{itemize}

\subsection{Registry Structure}

A GitHub repository with a \texttt{registry.json} index:

\begin{lstlisting}
{
  "extensions": [
    {
      "name": "sankey",
      "version": "1.0.0",
      "author": "maya-torres",
      "repo": "github.com/maya-torres/lux-ext-sankey",
      "ref": "v1.0.0",
      "render_mode": "texture",
      "rating": 4.2,
      "downloads": 347,
      "verified": {
        "lint": true,
        "ci": true,
        "tests": 12,
        "coverage": 0.85
      }
    }
  ]
}
\end{lstlisting}

\subsection{CLI}

\begin{center}
\begin{tabular}{@{}lp{8cm}@{}}
\toprule
\textbf{Command} & \textbf{Purpose} \\
\midrule
\texttt{lux ext list}      & List installed extensions. \\
\texttt{lux ext search <q>} & Search community registry. \\
\texttt{lux ext install <name>} & Install from registry. \\
\texttt{lux ext remove <name>}  & Remove. \\
\texttt{lux ext publish}   & Publish to registry (requires account). \\
\texttt{lux ext info <name>} & Show metadata, ratings, verification. \\
\bottomrule
\end{tabular}
\end{center}

\subsection{Verification Badges}

Each community extension displays verification status, similar to
GitHub repository badges:

\begin{itemize}[nosep]
  \item \textbf{Lint passing} --- \texttt{ruff check} clean
  \item \textbf{CI green} --- tests pass on latest Lux version
  \item \textbf{Test coverage} --- percentage reported
  \item \textbf{Community rating} --- 1--5 stars from users
  \item \textbf{Download count} --- adoption signal
  \item \textbf{Author verified} --- GitHub identity confirmed
\end{itemize}

This allows consumers to assess extension quality without reading
the source code.  The producer/consumer asymmetry is key: a small
number of authors (who enable LLM authoring or write extensions
manually) produce extensions consumed by a much larger number of
users who never need to author code.

% ===================================================================
\section{Frame Management}
% ===================================================================

Inherited from v1 with no structural changes.  Frames are the
organizing primitive:

\begin{itemize}[nosep]
  \item Created on first scene with a \field{frame\_id}.
  \item Shared ownership via \field{owner\_fds} set.
  \item Orphan model: disconnect $\to$ scenes persist with
    \texttt{\_ORPHAN\_FD} sentinel.
  \item Dock bar for minimized frames.
  \item Layout modes: tab (default) and stack.
  \item Frame flags: \field{no\_resize}, \field{no\_collapse},
    \field{no\_title\_bar}, \field{no\_background},
    \field{no\_scrollbar}, \field{auto\_resize}.
\end{itemize}

The service API adds introspection (\api{GET /api/frames}) and
remote frame creation (\api{POST /api/show} with
\field{frame\_id}).

% ===================================================================
\section{Security Model}
% ===================================================================

\subsection{Trust Boundaries}

\begin{center}
\begin{tabular}{@{}lp{4cm}p{5cm}@{}}
\toprule
\textbf{Boundary} & \textbf{Threat} & \textbf{Mitigation} \\
\midrule
Socket/HTTP & Unauthorized connection & Localhost binding by
  default.  Token auth for TCP/TLS. \\
Extension install & Malicious code & Consent dialog shows source.
  Community verification badges. \\
Extension render & Crash or hang & Exception isolation.  Auto-disable
  after repeated crashes. \\
RunPython command & Arbitrary exec & Consent-gated per permission
  mode.  Same trust as Bash tool. \\
Community registry & Supply chain & Verification badges (lint, CI,
  tests).  Pinned refs (git SHA). \\
\bottomrule
\end{tabular}
\end{center}

\subsection{Consent Model}

Extension installation follows the same consent pattern as v1's
render functions:

\begin{enumerate}[nosep]
  \item User or agent requests installation.
  \item Lux shows the extension source code in a consent dialog.
  \item User clicks Allow or Deny.
  \item Allowed: extension installs and registers.
  \item Denied: extension is not loaded.  No partial state.
\end{enumerate}

Community extensions with high ratings and green verification
badges may be auto-approved in Accept-edits or Bypass permission
modes.

% ===================================================================
\section{Migration from v1}
% ===================================================================

\begin{enumerate}
  \item \textbf{Phase 1: Registry.}  Refactor element dispatch from
    hardcoded \texttt{if/elif} chain to registry lookup.  All 27
    existing kinds become built-in extensions loaded at startup.
    Wire protocol unchanged.  MCP tools unchanged.  Zero user-facing
    change.
  \item \textbf{Phase 2: Service API.}  Add HTTP server alongside
    existing Unix socket.  MCP adapter calls service API internally.
    Both interfaces active.
  \item \textbf{Phase 3: User extensions.}  Enable
    \texttt{\char`\~/.lux/extensions/} loading.  Ship reference
    extensions.  \texttt{lux ext} CLI.
  \item \textbf{Phase 4: Agent SDK.}  Add LuxExchange, commands,
    permission policy.  Optional --- no impact on users who do not
    enable it.
  \item \textbf{Phase 5: Community.}  Launch registry.  Publish
    starter pack + reference extensions.  Verification pipeline.
\end{enumerate}

Each phase is independently shippable.  Phase 1 alone improves
the codebase (removes hardcoded dispatch) without changing the
product.  Phase 3 is where Lux~v2 becomes a different product
from v1.

\end{document}
